\title{The Era of 1-bit LLMs: All Large Language Models are in 1.58 Bits}
This version incorporates several changes relative to the original BitNet design, which are outlined as follows. Nevertheless, the exponential growth in LLM size has brought deployment difficulties to the forefront and intensified worries regarding their sustainability and operational costs, primarily due to substantial power usage. Expanding SRAM to alleviate this is not ideal, as it tends to be far more expensive than DRAM. Latency and memory measurements utilized a 2-bit kernel implementation, indicating further opportunities exist to optimize and minimize resource usage. There is further opportunity to compress these activations losslessly to 4 bits, or even lower, when working with 1.58-bit LLMs—an avenue to be explored in subsequent studies. Our focus was on matrix multiplication, as this operation constitutes the most substantial computational cost within LLMs. \textbf{\text{BitNet b1.58} establishes a new paradigm in scaling model efficiency with respect to both performance and inference overhead}. This approach omits zero-point quantization, thereby simplifying both the software implementation and hardware optimization. \item The 70B BitNet b1.58 outperforms the 13B FP16 LLM across latency, memory usage, and energy usage metrics. The results demonstrate that \text{BitNet b1.58} consistently outperforms across all evaluated tasks, underscoring the robust generalization capability of 1.58-bit LLMs. Consistent with perplexity analysis, task-level results indicate that \text{BitNet b1.58} at 3.9B exceeds the accuracy of \text{LLaMA LLM} 3B while operating at reduced memory and latency costs. The approach involves rescaling the weight matrix by its mean absolute value, and subsequently mapping each element to the closest integer within the permitted set:

\begin{equation}
    \widetilde{W} = \mathrm{RoundClip}(\frac{W}{\gamma + \epsilon}, -1, 1),
\end{equation}

\begin{equation}
    \mathrm{RoundClip}(x, a, b) = \max(a, \min(b, \mathrm{round}(x))),
\end{equation}

\begin{equation}
    \gamma = \frac{1}{nm}\sum_{ij} |W_{ij}|. Furthermore, the 3.9B version of \text{BitNet b1.58} demonstrates a 2.4-fold speed increase and 3.32-fold lower memory usage compared to \text{LLaMA LLM} 3B, while also surpassing its performance. A primary obstacle for efficient long-sequence inference is the high memory usage driven by key-value (KV) caches. Table~\ref{tab:ppl} provides an overview of perplexity outcomes and computational costs for both \text{BitNet b1.58} and \text{LLaMA LLM}. These advances have set new benchmarks across diverse natural language understanding and generation applications. The data suggest that as model capacity grows, the performance disparity between \text{BitNet b1.58} and \text{LLaMA LLM} becomes less pronounced. Relying on the energy estimation models presented in~\cite{energycost, pokebnn}, \text{BitNet b1.58} demonstrates a 71.4-fold reduction in energy consumption for matrix multiplication operations on 7nm process technology. By contrast, BitNet leverages integer addition for matrix multiplications, which substantially decreases the energy consumed during LLM operations. As depicted in Figure~\ref{fig:latency-memory-energy}, both latency and memory usage show clear scaling trends; notably, the acceleration effect is amplified as the models become larger. These findings indicate that \text{BitNet b1.58} represents a Pareto-efficient advance relative to current large language model architectures. The GPU memory consumption and inference latency of both \text{LLaMA LLM} and \text{BitNet b1.58} were quantitatively measured. 
\begin{abstract}
Building upon advances like BitNet~\cite{bitnet}, the landscape of Large Language Models (LLMs) is shifting towards highly efficient 1-bit representations. Details of zero-shot accuracy across downstream tasks can be found in Table~\ref{tab:zero-shot}. Crucially, \text{BitNet b1.58} matches the full-precision baseline’s results starting at the 3B scale. These benchmarks utilized the FasterTransformer\footnote{\url{https://github.com/NVIDIA/FasterTransformer}} library, which is optimized for low-latency inference on GPUs. This contributes to overall gains in inference speed and efficiency. \textbf{New Hardware for 1-bit LLMs}

Developments such as~Groq\footnote{\url{https://groq.com/}} have showcased the potential of purpose-built hardware (e.g., LPUs) for supporting LLM workloads. This design choice facilitates straightforward deployment within popular open-source platforms such as Huggingface, vLLM~\cite{vllm}, and llama.cpp\footnote{\url{https://github.com/ggerganov/llama.cpp}}, requiring minimal integration effort. Memory demands also scale in a comparable manner, primarily because the full-precision embedding layer's contribution to total memory footprint becomes proportionally less significant in larger models. Additionally, validation perplexity measurements were reported for the WikiText2~\cite{wikitext2} and C4~\cite{c4} benchmarks. These bottlenecks can be alleviated with 1.58-bit LLMs. Such limitations often prevent effective LLM deployment on these devices. StableLM 3b’s performance at 2T tokens is quoted directly from its official report. Additionally, we measured total energy usage for generating outputs with 512 tokens. Despite their prevalence in the deployment of commercial LLMs, post-training quantization generally yields sub-optimal results compared to alternatives. Additionally, transmission of activations over the network becomes considerably less burdensome. \section{Discussion and Future Work}

\textbf{1-bit Mixture-of-Experts (MoE) LLMs}

Mixture-of-Experts (MoE) architectures offer a resource-efficient solution for scaling LLMs, drastically cutting required FLOPs. \paragraph{Energy}

We analyzed the energy requirements attributable to arithmetic operations for both \text{BitNet b1.58} and \text{LLaMA LLM}. Our empirical results indicate that this change has an insignificant impact on model performance. \item The 30B BitNet b1.58 model exceeds the 7B FP16 LLM in latency, memory efficiency, and energy consumption. \paragraph{Quantization Function.} In order to restrict weight values to the set $\{-1, 0, +1\}$, we implement an \emph{absmean} based quantization procedure. For example, at 70B parameters, \text{BitNet b1.58} achieves a speed-up of 4.1x relative to the \text{LLaMA LLM} baseline. \end{itemize}

\paragraph{Training with 2T Tokens}

The scale of training tokens used plays a fundamental role in the effectiveness of LLMs. This progression has moved from standard 16-bit precision toward much lower bit-widths, including prominent 4-bit models~\cite{gptq, awq}. As reported in Table~\ref{tab:throughput}, \text{BitNet b1.58} 70B is capable of sustaining a batch size up to 11 times larger than that of \text{LLaMA LLM}, thereby achieving a throughput that is 8.9 times greater. At the same time, it brings considerable advantages in terms of inference latency, memory utilization, computational throughput, and energy efficiency. \end{abstract}

\section{The Era of 1-bit LLMs}

The last few years have witnessed extraordinary expansion in the scale and proficiency of Large Language Models (LLMs) within the AI domain. \paragraph{LLaMA-alike Components.}

The LLaMA~\cite{llama, llama2} architecture serves as the prevailing foundation among open-source large language models. The second benefit, as demonstrated in our experiments, is that \text{BitNet b1.58} can achieve parity with full-precision FP16 models regarding perplexity and downstream task performance, for models at least 3B in size, given identical model and training configurations. To maintain compatibility with this ecosystem, we structure \text{BitNet b1.58} with several LLaMA-inspired elements: employing RMSNorm~\cite{rmsnorm}, integrating SwiGLU~\cite{swiglu}, using rotary embedding~\cite{rope}, and eliminating all biases from the model. We documented the time per generated output token, recognizing this as the main contributor to inference computational cost. Beyond computation, inference-time parameter transfer—from DRAM to on-chip accelerator memory (like SRAM)—also incurs notable costs. In addition, 1.58-bit LLMs exhibit strong compatibility with CPUs, which predominate in these systems, enabling efficient execution of \text{BitNet b1.58} and further augmenting device functionality and performance. Here, we present a notable 1-bit LLM variant, designated \textbf{\text{BitNet b1.58}}, wherein all parameters are ternary, capable of taking the values \{-1, 0, 1\}. \textbf{LLMs on Edge and Mobile}

Leveraging 1.58-bit LLMs can bring substantial advantages for language models operating on edge and mobile platforms, which are frequently constrained by memory capacity and computational throughput. Both models were subsequently evaluated on a benchmark suite comprising Winogrande~\cite{winoGrande}, PIQA~\cite{piqa}, SciQ~\cite{sciq}, LAMBADA~\cite{lambada}, and ARC-easy~\cite{arc}. \section{BitNet b1.58}

\text{BitNet b1.58} builds upon the BitNet framework, a variant of the Transformer architecture wherein the conventional \emph{nn.Linear} layers are substituted with \emph{BitLinear} layers. \text{BitNet b1.58} inherits the central computational innovations of its predecessor, nearly eliminating multiplication in matrix operations and enabling highly optimized computation. However, deployment is often hindered by substantial memory demands and the cost of inter-chip data communication. \text{BitNet b1.58} makes significant advances in this regard by reducing activation precision from 16 to 8 bits, thereby enabling twice the sequence length with the same hardware constraints. Moreover, \text{BitNet b1.58} provides two further benefits. Recent advancements in 1-bit model architectures—exemplified by BitNet~\cite{bitnet}—suggest a promising pathway for cutting down large language model (LLM) operational costs without compromising accuracy. Our findings indicate that as model size increases, \text{BitNet b1.58} achieves increasingly greater energy efficiency over the FP16-based \text{LLaMA LLM}, owing to the expanding dominance of \emph{nn.Linear} operations in large models and the relatively smaller overhead from other architectural components. This improvement arises because the execution time for \emph{nn.Linear} layers increases with greater model size. The batch size was incrementally increased until the available GPU memory was fully utilized, maintaining a sequence length of 512. \end{equation}

For activation quantization, we utilize the same scheme as BitNet. Despite this drastic reduction in precision, our ternary LLM achieves parity with its full-precision Transformer counterpart (using FP16 or BF16) in both perplexity and downstream task accuracy for equivalent model sizes and training data volumes. Ideally, this would eliminate interconnect overhead if the complete model fits onto a single device. Drawing from the observations in Figures~\ref{fig:latency-memory-energy} and~\ref{fig:energy}, we present the following approximate correspondences for model sizes between the 1.58-bit and 16-bit representations:

\begin{itemize}
    \item A 13B BitNet b1.58 model surpasses the 3B FP16 LLM regarding latency, memory footprint, and energy expenditure. The additional zero augments the original binary BitNet, resulting in a representation size of 1.58 bits per parameter in binary terms. Compared to full-precision alternatives, 1-bit LLMs achieve significant reductions in both memory capacity and bandwidth requirements, ultimately lowering both the financial and temporal costs of model weight transfers from DRAM. To assess the ability of \text{BitNet b1.58} to scale with the number of tokens, we conducted training on a \text{BitNet b1.58} model with 2T tokens, adhering to the data preparation methodology of StableLM-3B~\cite{StableLM-3B-4E1T}, recognized as a leading open-source 3B parameter model. \paragraph{Memory and Latency}

We extended our experiments by increasing the model sizes to 7B, 13B, and 70B, and examined the associated computational cost. This approach further unlocks a fresh computing paradigm and paves the way for the development of specialized hardware geared towards 1-bit LLM architectures. \section{Results}

We conducted a comparative analysis between \text{BitNet b1.58} and our implementation of FP16 \text{LLaMA LLM} across multiple model sizes. \paragraph{Throughput}

We evaluated and contrasted the throughput performance of \text{BitNet b1.58} and \text{LLaMA LLM} (70B parameters), utilizing two 80GB A100 GPUs configured with pipeline parallelism~\cite{gpipe}, which is necessary to accommodate the 70B \text{LLaMA LLM}. Building upon this trajectory, there is a clear imperative to design novel systems and hardware architectures tailored for 1-bit LLMs, capitalizing on the new computational frameworks introduced with BitNet~\cite{bitnet}. Given that power is a central bottleneck in the computational capabilities of many hardware platforms, the reduction in energy usage further allows for more rapid execution. Evaluation was carried out following the protocols in \emph{lm-evaluation-harness}\footnote{\url{https://github.com/EleutherAI/lm-evaluation-harness}}. Conventional LLMs rely on 16-bit floating-point representations (such as FP16 or BF16), and their computation is largely dominated by matrix multiplications. The first is its increased expressive power, attributable to the ability of weights to explicitly mask features via the zero value—this addition enhances the representational and generalization capacity of 1-bit LLMs. The lowered energy and memory requirements of 1.58-bit LLMs facilitate their operation in resource-restricted environments, opening avenues for novel applications previously deemed impractical. In this study, we present a 1-bit LLM variant named \textbf{\text{BitNet b1.58}}, characterized by constraining every model parameter to one of three possible values: \{-1, 0, 1\}. Notably, this 1.58-bit representation inspires a novel scaling law and provides a blueprint for the next generation of LLMs that combine high effectiveness with low resource demands. \textbf{Native Support of Long Sequence in LLMs}

Long-context handling is a pressing need in contemporary LLM research. By lowering the precision of both weights and activations, these methods can dramatically reduce memory footprints and computational loads. An emerging strategy to mitigate these obstacles involves post-training quantization methods, which tailor model representations to use fewer bits for inference~\cite{smoothquant, gptq, quip, quip_sharp}. We reported zero-shot accuracy results in Table~\ref{tab:2t}; where tasks employ both accuracy and normalized accuracy, we report their average. It preserves the energy-efficient profile of the original BitNet and remains highly advantageous over FP16-based LLMs in memory usage, computational throughput, and latency. Both models were pre-trained for 100 billion tokens using the RedPajama dataset~\cite{redpajama} to provide an equitable evaluation setting. The model is initialized and trained entirely anew, employing weights quantized to 1.58 bits and activations at 8 bits. Consequently, the main computational expenses stem from floating-point addition and multiplication. This significantly broadens the potential of edge and mobile devices to leverage LLMs in diverse scenarios. Their zero-shot capabilities were assessed on an array of language understanding tasks, including ARC-Easy~\cite{arc}, ARC-Challenge~\cite{arc}, Hellaswag~\cite{hellaswag}, Winogrande~\cite{winoGrande}, PIQA~\cite{piqa}, OpenbookQA~\cite{openbookqa}, and BoolQ~\cite{boolq}. Notably, at the 3B parameter scale, \text{BitNet b1.58} achieves comparable perplexity to its full-precision \text{LLaMA LLM} counterpart, but with a 2.71-fold acceleration and a 3.55-fold reduction in GPU memory usage. The diminished memory usage of such models means that fewer hardware units are necessary for MoE deployment. For \text{BitNet b1.58}, we also integrated the 2-bit kernel provided by Ladder~\cite{ladder}. Figure~\ref{fig:energy} breaks down the sources of energy expenditure: the majority of operations in \text{BitNet b1.58} involve INT8 additions, whereas in \text{LLaMA LLM} most computations are FP16 additions and multiplications. However, rather than normalizing the activations before applying non-linear functions to fit within $[0, Q_b]$, we rescale all activations to the interval $[-Q_b, Q_b]$ on a per-token basis.